\begin{insight}{行列式测量的是有向面积}
把二阶矩阵的两列看成从原点出发的向量，行列式的绝对值就是它们张成的平行四边形面积，符号记录方向。倍加一列等于沿底边剪切图形：底和高都不变，因此行列式不变。交换两列则把有序方向反转，面积大小不变而符号改变。
\end{insight}
\begin{figure}[H]\centering
\begin{tikzpicture}[scale=.95]
\begin{scope}
\draw[laaxis](-.3,0)--(4,0) node[right]{$x$};\draw[laaxis](0,-.2)--(0,2) node[above]{$y$};
\filldraw[fill=Teal!12,draw=Teal](0,0)--(2,0)--(2.7,1.4)--(.7,1.4)--cycle;
\draw[lavec,Ink](0,0)--(2,0) node[midway,below]{$\vu$};
\draw[lavec,Gold](0,0)--(.7,1.4) node[above]{$\vv$};
\draw[dashed,Rust](.7,1.4)--(.7,0) node[midway,left]{$h$};
\node[lanote]at(1.7,-.65){$\det(\vu,\vv)=\text{底}\times\text{高}$};
\end{scope}
\begin{scope}[xshift=6cm]
\draw[laaxis](-.3,0)--(4,0) node[right]{$x$};\draw[laaxis](0,-.2)--(0,2) node[above]{$y$};
\filldraw[fill=Teal!12,draw=Teal](0,0)--(2,0)--(3.9,1.4)--(1.9,1.4)--cycle;
\draw[lavec,Ink](0,0)--(2,0) node[midway,below]{$\vu$};
\draw[lavec,Gold](0,0)--(1.9,1.4) node[above]{$\vv+t\vu$};
\draw[dashed,Rust](1.9,1.4)--(1.9,0) node[midway,right]{$h$};
\node[lanote]at(1.7,-.65){沿底边剪切，面积保持};
\end{scope}
\end{tikzpicture}
\caption{倍加列变换的面积解释。图中$t=0.6$，两图的底和高相同。}\label{fig:det-shear}
\end{figure}
\begin{clevertrick}{看到“另一列的倍数”，先消去它}
设$\vu=(2,0)\T,\vv=(1,3)\T$。按列线性性与重复列为零，
\[\det(\vu,\vv+7\vu)=\det(\vu,\vv)+7\det(\vu,\vu)=\det(\vu,\vv).\]
所以直接用$2\cdot3-0\cdot1=6$，不用先算出第二列$(15,3)\T$。若逐项计算，$2\cdot3-0\cdot15=6$也核验相同结果。几何上只是图\ref{fig:det-shear}中的剪切。若交换列次序，则
\[\det(\vv+7\vu,\vu)=-\det(\vu,\vv+7\vu)=-6,\]
不能只报面积6；有向面积还须保留次序给出的符号。
\end{clevertrick}
